\section{Evaluation}\label{sec:evaluation}

\newcommand\numbertodo[1]{{\color{red}{#1}}}
\def\exectotal{1176}
\def\spectotal{1279}
\def\lhtimetotal{148.76}
\def\lhsmttotal{1626\xspace}
\def\lhproofstotal{1524\xspace}
\def\pbetimetotal{150.88}
\def\pbesmttotal{4068\xspace}
\def\pbeproofstotal{638\xspace}

\def\lhtotaltotal{843}
\def\pbetotaltotal{306}

\newcolumntype{?}{!{\vrule width 1pt}}
\setlength{\tabcolsep}{3pt}

\begin{small}
\begin{table*}

\begin{center}
  \begin{tabular}{ ? l  ? r|r ? r|r|r ? r|r|r? }
    \hline
        \multirow{2}{*}{\textbf{Benchmark}}  & \multicolumn{2}{c?}{\textbf{Common}} & \multicolumn{3}{c?}{\textbf{Without PLE Search}} & \multicolumn{3}{c?}{\textbf{With PLE Search}} \\\cline{2-9}
          & Impl~(l) & Spec~(l)
          & Proof~(l)& Time~(s) & SMT~(q)
          & Proof~(l) & Time~(s) & SMT~(q) \\
    \hline
     \multicolumn{9}{?l?}{\textbf{Arithmetic}}\\
    \hline
    Fibonacci
    & 7 & 10
    & 38& 2.74 & 129
    & 16& 1.92 & 79
    \\
    Ackermann
    & 20 & 73
    & 196& 5.40        & 566
    & 119& 13.80       & 846
    \\
    \hline
    \multicolumn{9}{?l?}{\textbf{Class Laws}  Fig~\ref{fig:laws}}\\
    \hline
    Monoid
    & 33 & 50
    & 109& 4.47 & 34
    & 33& 4.22 & 209
    \\
    Functor
    & 48 & 44
    & 93& 4.97 & 26
    & 14& 3.68 & 68
    \\
    Applicative
    & 62 & 110
    & 241& 12.00 & 69
    & 74& 10.00 & 1090
    \\
    Monad
    & 63 & 42
    & 122& 5.39           & 49
    & 39& 4.89           & 250
    \\
    \hline
     \multicolumn{9}{?l?}{\textbf{Higher-Order Properties}}\\
    \hline
    Logical Properties %% this is NatInduction + NaturalDeduction
    & 0 & 20 % 2+18
    & 33& 2.71 & 32
    & 33& 2.74 & 32
    \\
    \hline
    Fold Universal
    & 10 & 44
    & 43& 2.17 & 24
    & 14& 1.46 & 48
    \\
    \hline
     \multicolumn{9}{?l?}{\textbf{Functional Correctness}}\\
    \hline
    SAT-solver
    & 92 & 34
    & 0& 50.00 & 50
    & 0& 50.00 & 50
    \\
    Unification
    & 51 & 60
    & 85& 4.77 & 195
    & 21& 5.64 & 422
    \\
    \hline
     \multicolumn{9}{?l?}{\textbf{Deterministic Parallelism}}\\
    \hline
    Conc. Sets
    & 597 & 329
    & 339& 40.10  & 339
    & 229& 40.70  & 861
    \\
    $n$-body
    & 163 & 251
    & 101& 7.41 & 61
    & 21& 6.27 & 61
    \\
    Par. Reducers
    & 30 & 212
    & 124& 6.63 & 52
    & 25& 5.56 & 52
    \\
     \hline
    \textbf{Total}
    & \exectotal   & \spectotal
    & \lhproofstotal& \lhtimetotal & \lhsmttotal
    & \pbeproofstotal& \pbetimetotal & \pbesmttotal
    \\
    \hline
	\end{tabular}
\end{center}
  \caption{We report verification \textbf{Time}
           (in seconds, on a 2.3GHz Intel{\textregistered} Xeon{\textregistered}
           CPU E5-2699 v3 with 18 physical cores and 64GiB RAM.),
           the number of \textbf{SMT} queries and size of
           \textbf{Proofs} (in lines).
           %
           The \textbf{Common} columns show sizes of common
           \textbf{Implementations} and \textbf{Specifications} .
           %
           We separately consider proofs \textbf{Without} and
           \textbf{With PLE Search}.
           %
           % \toolname (\lhtotaltotal~LoC in total) and
           % \toolname with \pbesym (\pbetotaltotal~LoC in total) into
           % \textbf{specifications}, \textbf{proof terms} and \textbf{executable} code.
         }
\label{table:evaluation}
\end{table*}
\end{small}

%
We have implemented reflection and \pbesym
in \toolname~\cite{Vazou14}.
%
Table~\ref{table:evaluation} summarizes our
evaluation which aims to determine
%
(1)~the kinds of programs and properties that can be verified,
%
(2)~how \pbesym simplifies \emph{writing} proofs, and
%
(3)~how \pbesym affects the verification time.


\mypara{Benchmarks}
%
% We used our approach to implement and verify
% a wide variety of programs.
%
We summarize our benchmarks below,
see \citep{appendix} for details.
%
\begin{itemize}[leftmargin=*]
%
\item\textbf{Arithmetic}
We proved arithmetic properties for the textbook
Fibonacci function (\cf~\S~\ref{sec:overview})
and the 12 properties of the Ackermann function
from~\cite{ackermann}.
%
\item\textbf{Class Laws}
%
We proved the monoid laws for the "Peano",
"Maybe" and "List" data types and the Functor,
Applicative, and Monad laws, summarized
in~Figure~\ref{fig:laws}, for the "Maybe",
"List" and "Identity" monads.
%
\item\textbf{Higher Order Properties}
%
We used natural deduction to prove textbook
logical properties as in~\S~\ref{sec:higher-order}.
%
% and this http://hume.ucdavis.edu/mattey/phi112/112dedurles_ho.pdf
%
We combined natural deduction principles with
\pbesym-search to
prove universality of right-folds, as
described in~\citep{foldrHutton} and
formalized in \agda~\citep{agdaequational}.
%
\item\textbf{Functional Correctness}
%
We proved correctness of a SAT solver
and a unification algorithm as implemented
in Zombie~\citep{Zombie}.
%
We proved that the SAT solver takes as
input a formula "f" and either returns "Nothing"
or an assignment that satisfies "f",
by reflecting the notion of satisfaction.
%
Then, we proved that if the
unification "unify s t" of two terms
"s" and "t" returns a substitution
"su", then applying "su" to "s" and "t"
yields identical terms.
%
Note that, while the unification function can itself
diverge, and hence cannot be reflected, our method
allows terminating and diverging functions
to soundly coexist.

\item\textbf{Deterministic Parallelism}
%
Retrofitting verification
onto an existing language with a mature parallel run-time
allows us to create three deterministic
parallelism libraries that, for the first time,
verify implicit assumptions about associativity
and ordering that are critical for determinism 
(\cf~\cite{appendix} for extended description).
%
First, we proved that the \emph{ordering laws} hold for keys
inserted into LVar-style concurrent sets \cite{kuper2014freeze}.
%
Second, we used "monad-par" \cite{monad-par} to implement an $n$-body
simulation, whose correctness relied upon proving that a triple of "Real"
(implementing) 3-d acceleration was a "Monoid".
%
Third, we built a DPJ-style \cite{DPJ} parallel-reducers library
whose correctness relied upon verifying that the reduced
arguments form a "CommutativeMonoid", and which was the basis of
a parallel array sum.
\end{itemize}

\mypara{Proof Effort}
%
We split the total lines of code of our benchmarks
into three categories:
%
\textbf{Spec} represents the refinement types that
encode theorems, lemmas, and function \emph{specifications};
%
\textbf{Impl} represents the rest of the Haskell
code that defines executable functions;
%
\textbf{Proofs} represent the sizes of the Haskell proof terms
(\ie functions returning \typp).
%
Reflection and \pbesym are optionally
enabled using pragmas; the latter is enabled
either for a whole file/module or per
top-level function.

\mypara{Runtime Overhead}
%
Proof terms have \emph{no} runtime overhead
as they will never be evaluated.
%
When verification of an executable term
depends on theorems, we use the below 
"withTheorem" function
%
\begin{mcode}
    withTheorem :: x:a -> $\typp$ -> { v:a | v = x }
    withTheorem x _ = x  
\end{mcode}
%
that inserts the proof argument into the static 
verification environment, relying upon laziness
to not actually evaluate the proof.
%
For example, when verification depends on the associativity 
of append on the lists "xs", "ys", and "zs", 
the invocation "withTheorem xs (app_assoc xs ys zs)"
extends the (static) SMT verification environment
with the instantiation of the associativity theorem of 
Figure~\ref{fig:list-laws}. 
%
This invocation adds no runtime overhead, since 
even though "app_assoc xs ys zs" is an expensive 
recursive function, it will never actually get evaluated. 
%
To ensure that proof terms are not evaluated in runtime, 
without using laziness, one can add one rewrite rule for 
each proof term that replaces the term with unit.
%
For example, the rewrite rule for "app_assoc" is 
%
\begin{code}
    RULES "assoc/runtime"  forall xs ys zs. app_assoc xs ys zs = () 
\end{code}
%
Such rules are sound, since each proof term is total, 
thus provably reduces to unit.

\mypara{Results}
%
The highlights of our evaluation are the following.
%
(1)~Reflection allows for the specification and verification
    of a wide variety of important properties of programs.
%
(2)~\pbesym drastically reduces the proof effort: by a factor
    of $2-5\times$ --- shrinking the total lines of
    proof from \lhproofstotal to \pbeproofstotal ---
    making it quite modest, about the size of
    the specifications of the theorems.
    Since \pbesym searches for equational
    properties, there are some proofs, that
    rarely occur in practice, that
    \pbesym cannot simplify, \eg the
    logical properties
    from~\S~\ref{sec:higher-order}.
%
(3)~\pbesym does not impose a performance penalty: even though
    proof search can make an order of magnitude many more SMT
    queries --- increasing the total SMT queries from
    \lhsmttotal without \pbesym to \pbesmttotal with \pbesym---
    most of these queries are simple and it is typically
    \emph{faster} to type-check the compact proofs enabled by
    \pbesym than it is to type-check the $2-5\times$ longer
    explicit proofs written by a human.

\begin{figure}[t!]
\begin{center}
\small
\begin{tabular}{rlrl}
\toprule

\multicolumn{2}{c}{\textbf{Monoid} (for Peano, Maybe, List)}
&
\multicolumn{2}{c}{\textbf{Functor} (for Maybe, List, Id)} \\

{Left Id.}  & $\emempty\ x\ \emappend\  \equiv x$
&
{Id.}    & $\efmap\ \eid\ xs \equiv \eid\ xs$ \\

{Right Id.} & $x\ \emappend\ \emempty \equiv x$
&
{Distr.} & $\efmap\ (g\ecompose\ h)\ xs \equiv (\efmap\ g\ \ecompose\ \efmap\ h)\ xs$\\

{Assoc.}     & $(x\ \emappend\ y)\ \emappend\ z \equiv x\ \emappend\ (y\ \emappend\ z)$
&
& \\

\midrule

\multicolumn{2}{c}{\textbf{Applicative} (for Maybe, List, Id)}
&
\multicolumn{2}{c}{\textbf{Monad} (for Maybe, List, Id)} \\

{Id.}    & $\epure \eid \eseq\ v \equiv v$
&
{Left Id.}   & $\ereturn\ a \ebind f \equiv f\ a$ \\

{Comp.}  & $\epure (\ecompose) \eseq u \eseq v \eseq w \equiv u \eseq (v \eseq w)$
&
{Right Id.}  & $m \ebind \ereturn \equiv m$ \\

{Hom.}   & $\epure\ f\ \eseq\ \epure\ x \equiv \epure\ (f\ x)$
&
{Assoc.}         & $(m\ebind f) \ebind g \equiv m\ebind (\lambda x \rightarrow f\ x \ebind g)$\\

{Inter.} & $u\ \eseq\ \epure\ y \equiv \epure\ (\$\ y) \ \eseq \ u$
&
& \\

\midrule

\multicolumn{2}{c}{\textbf{Ord} (for Int, Double, Either, $(,)$)}
&
\multicolumn{2}{c}{\textbf{Commutative Monoid} (for Int, Double, $(,)$)} \\

{Refl.}    & $x \leq x$
&
{Comm.}   & $x\ \emappend\ y \equiv y\ \emappend\ x$ \\

{Antisym.} & $x \leq y \land y \leq x \implies x \equiv y$
\\
{Trans.}   & $x \leq y \land y \leq z \implies x \leq z$
&
& \emph{(including Monoid laws)}
\\
{Total.}   & $x \leq y \lor y \leq x$ \\
%
\bottomrule
\end{tabular}
\end{center}
\caption{Summary of Verified Typeclass Laws.}
\label{fig:laws}

\end{figure}

% The source counts were generated by:
% $ cd benchmarks/haskell16/pos/
% $ find . -type f -name '*.hs' -exec sed -i '' s/\{-"/\{-#LH/ {} +
% $ sloccount
% For verified-instances:
% $ find src/Data/VerifiedEq* src/Data/VerifiedOrd* -type f -name '*.hs' -exec sloccount {} \+
% For LVish:
% $ git checkout master ; sloccount src/lvish/Data/LVar/PureSet.hs src/lvish/Data/LVar/SLSet.hs
% $ git checkout verified ; sloccount src/lvish/Data/LVar/PureSet.hs src/lvish/Data/LVar/SLSet.hs
% $ sloccount allpairs.hs ; sloccount allpairs_verified.hs
% $ sloccount IntegerSumReduction2NoVerification.hs ; sloccount IntegerSumReduction2.hs



%% \begin{figure}[t!]
%% \begin{minipage}[b]{0.47\textwidth}
%% % find . -type f -name '*.hs' -exec sed -r -i 's/\{-@/\{-#LH/' {} \+
% sloccount
% \begin{center}
\setlength{\tabcolsep}{3pt}
\small
\begin{tabular}{lllr}
\toprule
  \multicolumn{3}{l}{\textbf{CATEGORY}}              & \textbf{LOC} \\
\toprule
  \textbf{I.} & \multicolumn{3}{l}{\textbf{Arithmetic}} \\[0.05in]
   & Fibonacci      & \S~\ref{sec:overview}          &  48 \\ % Overview.hs
   & Ackermann      & \citep{ackermann}              & 280 \\ % Ackermann.hs

  \midrule

  \textbf{II.} & \multicolumn{3}{l}{\textbf{Algebraic Data Types}} \\[0.05in]

  & Fold Universal & \citep{agdaequational}          & 105 \\ % FoldrUniversal.hs
%  & Fold Fusion    & \citep{agdaequational}          &     \\

  \midrule

  \textbf{III.} & \textbf{Typeclass Laws} & Fig~\ref{fig:laws} & \\[0.05in]
  & Monoid         & \tPeano, \tMaybe, \tList        & 189 \\ % Monoid*.hs
  & Functor        & \tMaybe, \tList, \tId, \tReader & 296 \\ % Functor*.hs - FunctorReader.hs
  & Applicative    & \tMaybe, \tList, \tId, \tReader & 578 \\ % Applicative*.hs
  & Monad          & \tMaybe, \tList, \tId, \tReader & 435 \\ % Monad*.hs

  \midrule

  \textbf{IV.} & \multicolumn{3}{l}{\textbf{Functional Correctness}} \\[0.05in]
  & SAT Solver     & \citep{Zombie}                  & 133 \\ % Solver.hs
  & Unification    & \citep{Sjoberg2015}             & 200 \\ % Unification

  \midrule

  \textbf{V.} & \multicolumn{3}{l}{\textbf{Deterministic Parallelism}} \\[0.05in]
  & Conc. Sets     & \S~\ref{sec:set}           & 1677 \\ % VerifiedEq/Ord + PureSet.hs + SLSet.hs
  & $n$-body       & \S~\ref{sec:nbody}         & 435  \\ % VerifiedCommutativeMonoid + Iso + allpairs_verified.hs
  & Par. Reducers  & \S~\ref{sec:reducer}       & 349  \\ % VerifiedCommutativeMonoid + Iso + IntegerSumReduction2.hs

  \midrule

  \multicolumn{3}{l}{\textbf{TOTAL}}                 & 4155 \\
\bottomrule
\end{tabular}
%\end{center}
\caption{\textbf{Summary of Case Studies}}
\label{fig:eval-summary}

%% \end{minipage}
%% \hspace{0.2in}
%% \begin{minipage}[b]{0.47\textwidth}
%% 
\begin{center}
\small
\begin{tabular}{rl}
\toprule
\multicolumn{2}{c}{\textbf{Monoid}} \\
{Left Id.}  & $\emempty\ x\ \emappend\  \equiv x$  \\
{Right Id.} & $x\ \emappend\ \emempty \equiv x$  \\
{Assoc}     & $(x\ \emappend\ y)\ \emappend\ z \equiv x\ \emappend\ (y\ \emappend\ z)$ \\

\midrule

\multicolumn{2}{c}{\textbf{Functor}} \\
{Id.}    & $\efmap\ \eid\ xs \equiv \eid\ xs$ \\
{Distr.} & $\efmap\ (g\ecompose\ h)\ xs \equiv (\efmap\ g\ \ecompose\ \efmap\ h)\ xs$\\

\midrule

\multicolumn{2}{c}{\textbf{Applicative}} \\

{Id.}    & $\epure \eid \eseq\ v \equiv v$ \\
{Comp.}  & $\epure (\ecompose) \eseq u \eseq v \eseq w \equiv u \eseq (v \eseq w)$ \\
{Hom.}   & $\epure\ f\ \eseq\ \epure\ x \equiv \epure\ (f\ x)$\\
{Inter.} & $u\ \eseq\ \epure\ y \equiv \epure\ (\$\ y) \ \eseq \ u$ \\

\midrule
\multicolumn{2}{c}{\textbf{Monad}} \\
{Left Id.}   & $\ereturn\ a \ebind f \equiv f\ a$ \\
{Right Id.}  & $m \ebind \ereturn \equiv m$ \\
{Assoc}         & $(m\ebind f) \ebind g \equiv m\ebind (\lambda x \rightarrow f\ x \ebind g)$\\
\bottomrule
\end{tabular}
\end{center}
\caption{\textbf{Summary of Verified Typeclass Laws}}
\label{fig:laws}

%% \end{minipage}
%% \end{figure}

\begin{comment}
We have implemented refinement reflection
in \toolname.
%
In this section, we evaluate our approach
by using \toolname to verify a variety of
deep specifications of Haskell functions
drawn from the literature and categorized
in Figure~\ref{fig:eval-summary},
totalling about 4,000 lines of specifications
and proofs.
%
Verification is fast: from 1 to 7s per
file (module), which is about twice as
much time as GHC takes to compile the
corresponding module.
%
Overall there is about one line of
type specification (theorems)
for three lines of code (implementations and proof).
%
Next, we detail each of the five classes of
specifications, illustrate how they were
verified using refinement reflection, and
discuss the strengths and weaknesses of
our approach.
%
\emph{All} of these proofs require refinement
reflection, \ie are beyond the scope of shallow
refinement typing.

%\RGS{Should we make a mention of VIFP here?}
%\NV{Proof strategies can go away}
\mypara{Proof Strategies.}
%
Our proofs use three building blocks, that are seamlessly
connected via refinement typing:
%
\begin{itemize}
  \item \emphbf{Un/folding}
     definitions of a function "f" at
     arguments "e1...en", which due
     to refinement reflection, happens
     whenever the term "f e1 ... en"
     appears in a proof.
     For exposition, we render the function
     whose un/folding is relevant as "#f#";
%
  \item \emphbf{Lemma Application}
     which is carried out by using
     the ``because'' combinator
     ($\because$) to instantiate
     some fact at some inputs;
%
  \item \emphbf{SMT Reasoning}
     in particular, \emph{arithmetic},
     \emph{ordering} and \emph{congruence closure}
     which kicks in automatically (and predictably!),
     allowing us to simplify proofs by not
     having to specify, \eg which subterms
     to rewrite.
\end{itemize}

\subsection{Arithmetic Properties} \label{subsec:arith} \label{subsec:ackermann}



\subsection{Algebraic Data Properties}
\label{subsec:fold}

%The second category of properties pertains to data types.
\mypara{Fold Universality}
%Next, we proved properties of list folding,
as encoded in \agda~\citep{agdaequational}.
%
The following encodes the universal property of "foldr"
%
\begin{code}
  foldr_univ :: f:(a -> b -> b)
             -> h:(L a -> b)
             -> e:b
             -> ys:L a
             -> base:{h [] == e }
             -> step:(x:a -> xs:[a] -> {h (x:xs) == f x (h xs)})
             -> {h ys == foldr f e ys}
\end{code}
%
The \toolname proof term of "foldr_univ"
is similar to the one in \agda.
But, there are two major differences.
%
First, specifications do not support
quantifiers: universal quantification
of "x" and "xs" in the step assumption
is encoded as functional arguments.
%
Second, unlike \agda, \toolname does not
support optional arguments, thus to use
"foldr_univ" the proof arguments "base"
and "step" that were implicit arguments
in \agda need to be explicitly provided.
%
For example, the following code calls
"foldr_universal" to prove "foldr_fusion"
by explicitly applying proofs for the
base and the step arguments
%
\begin{code}
  foldr_fusion :: h:(b -> c)
               -> f:(a -> b -> b)
               -> g:(a -> c -> c)
               -> e:b
               -> z:[a]
               -> x:a -> y:b
               -> fuse: {h (f x y) == g x (h y)})
               -> {(h.foldr f e) z == foldr g (h e) z}

  foldr_fusion h f g e ys fuse = foldr_univ g (h . foldr f e) (h e) ys
                                   (fuse_base h f e)
                                   (fuse_step h f e g fuse)
\end{code}
%
Where, for example, "fuse_base" is a function with type
%
\begin{code}
  fuse_base :: h:(b -> c) -> f:(a -> b -> b) -> e:b -> {(h . foldr f e) [] == h e}
\end{code}



\subsection{Typeclass Laws}

We used \toolname to prove the Monoid, Functor,
Applicative, and Monad Laws, summarized in
Figure~\ref{fig:laws}, for various user-defined
instances summarized in Figure~\ref{fig:eval-summary}.

%% The purpose of these proofs is to investigate the
%% proving abilities of \libname.
%% For this purpose, we defined the appropriate class
%% operators on user defined lists, instead of using
%% Haskell's predefined class instances.
%% %
%% In the near future, we plan to embed these proofs
%% to check the laws on real Haskell instances,
%% but this requires some engineering from the
%% \liquidHaskell team.

\mypara{Monoid Laws}
%
A Monoid is a datatype equipped with an associative
binary operator $\emappend$ ("mappend") and an \emph{identity}
element $\emempty$.
%
We prove that
%
"Peano" (with a suitable "add" and "Z"),
"Maybe" (with a suitable "mappend" and "Nothing"), and
"List" (with append "++" and "[]") satisfy the monoid laws.
%
%%For example, we prove that "++" (\S~\ref{subsec:list})
%%is associative by reifying the textbook proof~\cite{HuttonBook}
%%into a Haskell function, where the induction
%%corresponds to case-splitting and recurring
%%on the first argument:
%%%
%%\begin{mcode}
%%assoc :: xs:[a] -> ys:[a] -> zs:[a] ->
%%       {(xs ++ ys) ++ zs = xs ++ (ys ++ zs)}
%%
%%assoc [] ys zs     = ([] ++ ys) ++ zs
%%                   =. [] ++ (ys ++ zs)
%%                   ** QED
%%assoc (x:xs) ys zs = ((x:  xs)++ ys) ++ zs
%%                   =. (x: (xs ++ ys)) ++ zs
%%                   =.  x:((xs ++ ys) ++ zs)
%%                   =.  x: (xs ++ (ys ++ zs))
%%                       $\because$ assoc xs ys zs
%%                   =. (x:xs)  ++ (ys ++ zs)
%%                   ** QED
%%\end{mcode}
%%

\mypara{Functor Laws}
%
A type is a functor if it has a function
"fmap" that satisfies the \emph{identity}
and \emph{distribution} (or fusion) laws
in Figure~\ref{fig:laws}.
%
For example, consider the proof of
the "map" distribution law for the lists,
also known as ``map-fusion'', which is the
basis for important optimizations in
GHC~\cite{ghc-map-fusion}.
%
We reflect the definition of "map" for lists
%
%%\begin{code}
%%  reflect map :: (a -> b) -> [a] -> [b]
%%  map f []     = []
%%  map f (x:xs) = f x : fmap f xs
%%\end{code}
%
and use it to specify fusion and verify it by an inductive proof:
% by induction (recursion) on the list argument:
%
\begin{mcode}
  map_fusion :: f:(b -> c) -> g:(a -> b) -> xs:[a] -> {map (f . g) xs = (map f . map g) xs}
\end{mcode}

%%\begin{mcode}
%%  map_fusion f g []
%%    =  ((map f) #.# (map g)) []
%%    =. (map f) (#map# g [])
%%    =. #map# f []
%%    =. []
%%    =. #map# (f . g) []
%%    ** QED
%%
%%  map_fusion f g (x:xs)
%%    =   #map# (f . g) (x:xs)
%%    =. (f . g) x : map (f . g) xs
%%    =. (f . g) x : (map f #.# map g) xs
%%       $\because$  map_fusion f g xs
%%    =. (f# . #g) x : map f (map g xs)
%%    =. f   (g  x): map f (map g xs)
%%    =. #map# f (g x: map g xs)
%%    =. map f   (#map# g (x:xs))
%%    =. (map f #.# map g)(x:xs)
%%    ** QED
%%\end{mcode}
%%
% \NV{Say why we need defunctionalization}
% \NV{We use app function like HALO (link to the theory)}

% \mypara{Applicative}
\mypara{Monad Laws}
%
The monad laws, which relate the
properties of the two operators
$\ebind$ and $\ereturn$ (Figure~\ref{fig:laws}),
refer to $\lambda$-functions which
are encoded in our logic as uninterpreted functions.
%%thus their proof exercises
%%our support for defunctionalization
%%and $\eta$- and $\beta$-equivalence.
% and the extensionality axioms to prove.
%
For example, we can encode the associativity property
as a refinement type alias:
%
\begin{code}
  type AssocLaw m f g = { m >>= f >>= g = m >>= (\x -> f x >>= g) }
\end{code}
%
and use it to prove that the list-bind is associative:
%
\begin{mcode}
  assoc :: m:[a] -> f:(a ->[b]) -> g:(b ->[c]) -> AssocLaw m f g
  assoc [] f g     =  [] #>>=# f >>= g
                   =. [] #>>=# g
                   =. []
                   =. [] #>>=# (\x -> f x >>= g) ** QED
  assoc (x:xs) f g =  (x:xs) #>>=# f  >>= g
                   =. (f x ++ xs >>= f) >>= g
                   =. (f x >>= g) ++ (xs >>= f >>= g) $\because$ bind_append (f x) (xs >>= f) g
                   =. (f x >>= g) ++ (xs >>= \y -> f y >>= g) $\because$ assoc xs f g
                   =. (\y -> f y >>= g) x ++ (xs >>= \y -> f y >>= g) $\because$ $\beta$eq f g x
                   =. (x:xs) #>>=# (\y -> f y >>= g) ** QED
\end{mcode}
%
Where the "bind_append" fusion lemma states that:
%
\begin{code}
bind_append
  :: xs:[a] -> ys:[a] -> f:(a -> [b]) -> {(xs++ys) >>= f = (xs >>= f)++(ys >>= f)}
\end{code}
%
Notice that the last step requires
$\beta$-equivalence on anonymous
functions, which we get by explicitly
inserting the redex in the logic,
via the following lemma with "trivial" proof
%
\begin{mcode}
    $\beta$eq :: f:_ -> g:_ -> x:_ -> {f x >>= g = (\y -> f y >>= g) x}
    $\beta$eq _ _ _ = trivial
\end{mcode}

\subsection{Functional Correctness} \label{subsec:programs}

Next, we proved correctness of two programs
from the literature: a unification
algorithm and a SAT solver.

\mypara{Unification}
%
We verified the
unification of first order terms, as
presented in~\citep{Sjoberg2015}.
%
First, we define a predicate alias for
when two terms "s" and "t" are equal
under a substitution "su":
%
\begin{code}
  eq_sub su s t = apply su s == apply su t
\end{code}
%
Now, we defined a Haskell function
"unify s t" that can diverge, or return
"Nothing", or return a substitution "su"
that makes the terms equal:
%
\begin{code}
  unify :: s:Term -> t:Term -> Maybe {su | eq_sub su s t}
\end{code}
%
For specification and verification,
we only needed to reflect "apply" and
not "unify"; thus, we only had to verify
that the former terminates, and not the latter.
%
We proved correctness by invoking
separate helper lemmata.
%
For example, to prove the post-condition
when unifying a variable "TVar i" with
a term "t" in which "i" \emph{does not}
appear, we apply a lemma "not_in":
%
\begin{mcode}
  unify (TVar i) t2 | not (i Set_mem freeVars t2) = Just (const [(i, t2)] $\because$ not_in i t2)
\end{mcode}
%%
\ie if "i" is not free in "t",
the singleton substitution yields "t":
%
\begin{code}
  not_in :: i:Int -> t:{Term | not (i Set_mem freeVars t)} -> {eq_sub [(i,t)] (TVar i) t}
\end{code}
%
This example highlights the benefits of
partial verification on a legacy programming language:
potential diverging code (\eg the function "unify") coexists and invokes proof terms.

\mypara{SAT Solver}
%
As another example, we implemented and verified the simple
SAT solver used to illustrate and evaluate
the features of the dependently typed language
Zombie~\citep{Zombie}.
%
The solver takes as input a formula "f"
and returns an assignment that
\emph{satisfies} "f" if one exists,
as specified below.
%
\begin{code}
  solve :: f:Formula -> Maybe {a:Assignment | sat a f}
\end{code}
%
The function "sat a f" returns "True" iff the assignment
"a" satisfies the formula "f".
%
Verifying "solve" follows directly by
reflecting "sat" into the refinement logic.
\end{comment}
